% GENERATED FILE. Edit canonical lessons, then run npm run generate:books.
% canonical-source-hash: fnv1a64:7815cc78990f4e76
% canonical-lessons: ML-C477-katavu, ML-C477-pulayoram, ML-C477-kanal, ML-C477-nintikkatakkuka, ML-C477-kutikkuka

\chapter{By the River}
\label{ch:ml-by-the-river}
\hlchaptermodality{477}

\begin{chapteropening}
\textbf{By the end of this chapter:} \emph{I can say a landing place, a riverbank, a canal, to swim across and to dash forward.}

Complete the last lesson of chapter 477: I can say a landing place, a riverbank, a canal, to swim across and to dash forward.
\end{chapteropening}

\section[\texorpdfstring{\ml{കടവ്} (kaṭavŭ)}{കടവ് (kaṭavŭ)}]{\ml{കടവ്} (kaṭavŭ) — a landing place; a ghat}

\label{lesson:ML-C477-katavu}



\begin{quote}
Before the new one: say the Malayalam for to toss about, then the Malayalam for to come ashore.
\end{quote}

\subsection*{You'll want to know: \ml{കടവ്}}
\textbf{\ml{കടവ്}} — \emph{kaṭavŭ} — \textquotedblleft{}a landing place; a ghat\textquotedblright{}.

\begin{quote}
\ml{സ്ത്രീകൾ} \ml{രാവിലെ} \ml{കടവിൽ} \ml{തുണി} \ml{കഴുകി}.
\end{quote}

\begin{quote}
\emph{strīkaḷ rāvile kaṭavil tuṇi kaḻuki.}
\end{quote}

\begin{quote}
In the morning the women washed clothes at the landing.
\end{quote}

A spot on a river bank for bathing or for the ferry. Many place names end in -\ml{ക്കടവ്}.

\begin{cousinweb}
\textbf{\ml{കടവ്}} (\emph{kaṭavŭ}) is native Dravidian, from the root of \ml{കടക്കുക}, to cross: the place where you cross (DEDR 1109).
\end{cousinweb}

\begin{grammarlens}[title={What you've built}]
One more word for getting around.
\end{grammarlens}

\subsection*{Your turn}
\begin{itemize}
\raggedright
  \item \emph{Say it:} \emph{kaṭavŭ}
  \item \emph{Say it:} \emph{kaṭavŭ}, once more
  \item \emph{Say it:} \emph{strīkaḷ rāvile kaṭavil tuṇi kaḻuki.} — In the morning the women washed clothes at the landing.
  \item \emph{Recall:} say the Malayalam for to toss about, then the Malayalam for to come ashore, then say \emph{kaṭavŭ} again
\end{itemize}

\begin{tcolorbox}[breakable,colback=teal!4,colframe=teal!35!black,title={Before you move on}]
What does \ml{കടവ്} mean? (\textquotedblleft{}A landing place; a ghat\textquotedblright{}.) Say it once more.
\end{tcolorbox}
\section[\texorpdfstring{\ml{പുഴയോരം} (puḻayōraṁ)}{പുഴയോരം (puḻayōraṁ)}]{\ml{പുഴയോരം} (puḻayōraṁ) — a riverbank}

\label{lesson:ML-C477-pulayoram}



\begin{quote}
Before the new one: say the Malayalam for to come ashore, then the Malayalam for a landing place.
\end{quote}

\subsection*{You'll want to know: \ml{പുഴയോരം}}
\textbf{\ml{പുഴയോരം}} — \emph{puḻayōraṁ} — \textquotedblleft{}a riverbank\textquotedblright{}.

\begin{quote}
\ml{പുഴയോരത്ത്} \ml{ധാരാളം} \ml{തെങ്ങുകൾ} \ml{ഉണ്ട്}.
\end{quote}

\begin{quote}
\emph{puḻayōrattŭ dhārāḷaṁ teṅṅukaḷ uṇṭŭ.}
\end{quote}

\begin{quote}
There are many coconut palms on the riverbank.
\end{quote}

Built like \ml{വഴിയോരം}: \ml{പുഴ}, a river, plus \ml{ഓരം}, an edge.

\begin{cousinweb}
\textbf{\ml{പുഴയോരം}} (\emph{puḻayōraṁ}) is a native Dravidian compound of \ml{പുഴ}, a river (DEDR 4318), and \ml{ഓരം}, an edge (DEDR 1060).
\end{cousinweb}

\begin{grammarlens}[title={What you've built}]
One more word for getting around.
\end{grammarlens}

\subsection*{Your turn}
\begin{itemize}
\raggedright
  \item \emph{Say it:} \emph{puḻayōraṁ}
  \item \emph{Say it:} \emph{puḻayōraṁ}, once more
  \item \emph{Say it:} \emph{puḻayōrattŭ dhārāḷaṁ teṅṅukaḷ uṇṭŭ.} — There are many coconut palms on the riverbank.
  \item \emph{Recall:} say the Malayalam for to come ashore, then the Malayalam for a landing place, then say \emph{puḻayōraṁ} again
\end{itemize}

\begin{tcolorbox}[breakable,colback=teal!4,colframe=teal!35!black,title={Before you move on}]
What does \ml{പുഴയോരം} mean? (\textquotedblleft{}A riverbank\textquotedblright{}.) Say it once more.
\end{tcolorbox}
\section[\texorpdfstring{\ml{കനാൽ} (kanāl)}{കനാൽ (kanāl)}]{\ml{കനാൽ} (kanāl) — a canal}

\label{lesson:ML-C477-kanal}



\begin{quote}
Before the new one: say the Malayalam for a landing place, then the Malayalam for a riverbank.
\end{quote}

\subsection*{You'll want to know: \ml{കനാൽ}}
\textbf{\ml{കനാൽ}} — \emph{kanāl} — \textquotedblleft{}a canal\textquotedblright{}.

\begin{quote}
\ml{ഗ്രാമത്തിലൂടെ} \ml{ഒരു} \ml{ചെറിയ} \ml{കനാൽ} \ml{ഒഴുകുന്നു}.
\end{quote}

\begin{quote}
\emph{grāmattilūṭe oru ceṟiya kanāl oḻukunnu.}
\end{quote}

\begin{quote}
A small canal flows through the village.
\end{quote}

Canals join the backwaters and carry boats and water for the fields.

\begin{cousinweb}
\textbf{\ml{കനാൽ}} (\emph{kanāl}) was borrowed from English \emph{canal}, a channel dug for water or boats.
\end{cousinweb}

\begin{grammarlens}[title={What you've built}]
One more word for getting around.
\end{grammarlens}

\subsection*{Your turn}
\begin{itemize}
\raggedright
  \item \emph{Say it:} \emph{kanāl}
  \item \emph{Say it:} \emph{kanāl}, once more
  \item \emph{Say it:} \emph{grāmattilūṭe oru ceṟiya kanāl oḻukunnu.} — A small canal flows through the village.
  \item \emph{Recall:} say the Malayalam for a landing place, then the Malayalam for a riverbank, then say \emph{kanāl} again
\end{itemize}

\begin{tcolorbox}[breakable,colback=teal!4,colframe=teal!35!black,title={Before you move on}]
What does \ml{കനാൽ} mean? (\textquotedblleft{}A canal\textquotedblright{}.) Say it once more.
\end{tcolorbox}
\section[\texorpdfstring{\ml{നീന്തിക്കടക്കുക} (nīntikkaṭakkuka)}{നീന്തിക്കടക്കുക (nīntikkaṭakkuka)}]{\ml{നീന്തിക്കടക്കുക} (nīntikkaṭakkuka) — to swim across}

\label{lesson:ML-C477-nintikkatakkuka}



\begin{quote}
Before the new one: say the Malayalam for a riverbank, then the Malayalam for a canal.
\end{quote}

\subsection*{You'll want to know: \ml{നീന്തിക്കടക്കുക}}
\textbf{\ml{നീന്തിക്കടക്കുക}} — \emph{nīntikkaṭakkuka} — \textquotedblleft{}to swim across\textquotedblright{}.

\begin{quote}
\ml{ചെറുപ്പത്തിൽ} \ml{അച്ഛൻ} \ml{ആ} \ml{പുഴ} \ml{നീന്തിക്കടന്നു}.
\end{quote}

\begin{quote}
\emph{ceṟuppattil acchan ā puḻa nīntikkaṭannu.}
\end{quote}

\begin{quote}
When he was young, father swam across that river.
\end{quote}

Past: \ml{നീന്തിക്കടന്നു}, \emph{nīntikkaṭannu}.

\ml{നീന്തി}, having swum, plus \ml{കടക്കുക}, to cross. What is crossed takes no ending.

\begin{cousinweb}
\textbf{\ml{നീന്തിക്കടക്കുക}} (\emph{nīntikkaṭakkuka}) is a native Dravidian compound of \ml{നീന്തുക}, to swim (DEDR 3687), and \ml{കടക്കുക}, to cross (DEDR 1109).
\end{cousinweb}

\begin{grammarlens}[title={What you've built}]
One more word for getting around.
\end{grammarlens}

\subsection*{Your turn}
\begin{itemize}
\raggedright
  \item \emph{Say it:} \emph{nīntikkaṭakkuka}
  \item \emph{Say it:} \emph{nīntikkaṭakkuka}, once more
  \item \emph{Say it:} \emph{ceṟuppattil acchan ā puḻa nīntikkaṭannu.} — When he was young, father swam across that river.
  \item \emph{Recall:} say the Malayalam for a riverbank, then the Malayalam for a canal, then say \emph{nīntikkaṭakkuka} again
\end{itemize}

\begin{tcolorbox}[breakable,colback=teal!4,colframe=teal!35!black,title={Before you move on}]
What does \ml{നീന്തിക്കടക്കുക} mean? (\textquotedblleft{}To swim across\textquotedblright{}.) Say it once more.
\end{tcolorbox}
\section[\texorpdfstring{\ml{കുതിക്കുക} (kutikkuka)}{കുതിക്കുക (kutikkuka)}]{\ml{കുതിക്കുക} (kutikkuka) — to dash forward; to surge}

\label{lesson:ML-C477-kutikkuka}



\begin{quote}
Before the new one: say the Malayalam for a canal, then the Malayalam for to swim across.
\end{quote}

\subsection*{You'll want to know: \ml{കുതിക്കുക}}
\textbf{\ml{കുതിക്കുക}} — \emph{kutikkuka} — \textquotedblleft{}to dash forward; to surge\textquotedblright{}.

\begin{quote}
\ml{പച്ച} \ml{സിഗ്നൽ} \ml{കണ്ടതും} \ml{വണ്ടികൾ} \ml{മുന്നോട്ട്} \ml{കുതിച്ചു}.
\end{quote}

\begin{quote}
\emph{pacca signal kaṇṭatuṁ vaṇṭikaḷ munnōṭṭŭ kuticcu.}
\end{quote}

\begin{quote}
As soon as they saw the green signal, the vehicles shot forward.
\end{quote}

Past: \ml{കുതിച്ചു}, \emph{kuticcu}.

Prices that shoot up also \ml{കുതിക്കുക}: \ml{വില} \ml{കുതിച്ചു}, the price shot up.

\begin{cousinweb}
\textbf{\ml{കുതിക്കുക}} (\emph{kutikkuka}) is native Dravidian, \textquotedblleft{}to jump, to leap\textquotedblright{} (DEDR 1705).
\end{cousinweb}

\begin{grammarlens}[title={What you've built}]
One more word for getting around.
\end{grammarlens}

\subsection*{Your turn}
\begin{itemize}
\raggedright
  \item \emph{Say it:} \emph{kutikkuka}
  \item \emph{Say it:} \emph{kutikkuka}, once more
  \item \emph{Say it:} \emph{pacca signal kaṇṭatuṁ vaṇṭikaḷ munnōṭṭŭ kuticcu.} — As soon as they saw the green signal, the vehicles shot forward.
  \item \emph{Recall:} say the Malayalam for a canal, then the Malayalam for to swim across, then say \emph{kutikkuka} again
\end{itemize}

\begin{tcolorbox}[breakable,colback=teal!4,colframe=teal!35!black,title={Before you move on}]
What does \ml{കുതിക്കുക} mean? (\textquotedblleft{}To dash forward; to surge\textquotedblright{}.) Say it once more.
\end{tcolorbox}
